\subsection{Motivation}
% \Khang{Nếu đổi thành đề tài Information Retrieval $\rightarrow$ Sửa lại cho phù hợp}

% \Khang{Cần thiết fashion nữa không? Có thể cần nhưng phải chú trọng feature extraction và multimodal attention}
% % In the evolving landscape of economics and retail, the fashion industry is experiencing rapid growth, catering to diverse demographics and preferences. However, as the number of fashion stores increases, consumers often find themselves spending considerable time searching for a store that aligns with their unique interests, financial constraints, and other specific requirements. This search predominantly relies on human-assisted interactions through websites, where customers engage in conversations with store managers.  To address this challenge, the development of a multimodal chatbot emerges as a potential solution. By integrating image and text into the conversation, the chatbot can efficiently assist customers in navigating through the diverse offerings and overall product information of the fashion stores. 

% \noindent As the fashion market expands to cater to diverse demographics and preferences, the conventional methods of searching for products become increasingly inefficient. Consumers often face challenges in finding fashion items that align with their specific preferences, budget constraints, and other unique requirements. Traditional approaches rely heavily on human-assisted interactions through websites, where customers engage in conversations with store managers to seek guidance or recommendations. However, this manual process is time-consuming, prone to errors, and fails to leverage the full potential of technology in enhancing customer experiences. The inefficiency of this process raises a significant challenge: How can this whole process be automated to enhance customer experience and streamline operations for fashion stores?

% \noindent By introducing a multimodal searching method, we aim to revolutionize the way customers interact with fashion stores online. The integration of image and text capabilities allows for a more intuitive and efficient communication channel between the consumer and the store. With this technology, customers can simply provide an image or description of the desired product, and the system can quickly analyze their preferences and recommend suitable options. Through fine-tuning a Vietnamese multimodal model and implementing an efficient information retrieval pipeline, the system can understand customer preferences, recommend appropriate products, and provide personalized assistance, thus revolutionizing the way customers interact with fashion stores online.

% \noindent Moreover, the development of such a system addresses the pressing need for automation in the fashion retail sector. As the number of fashion stores continues to rise, streamlining operations and improving customer experiences become paramount. A well-trained multimodal search system can effectively assist customers in navigating through the vast array of products, providing personalized recommendations, and answering queries regarding price, availability, and discounts.

% Paragraph 1: Bối cảnh, khó khăn hiện tại và đặt vấn đề (Tương đương đoạn 1 của bài mẫu)
In the evolving landscape of higher education and vocational training, the demand for job-ready graduates is experiencing rapid growth, catering to diverse industry requirements. However, as the gap between academic theory and practical application widens, students often find themselves spending considerable effort struggling to transition from theoretical learning to actual career readiness. This transition is further complicated by a modern cognitive crisis described as the "age of interruption," where the constant barrage of digital stimuli from smartphones and social media leads to "continuous partial attention". In this state, learners perpetually shift their focus across multiple stimuli without fully immersing themselves in any single task, resulting in a superficial understanding of information. Research indicates that students in the presence of technological distractors can often only maintain sustained focus for a mere six minutes before being diverted. Consequently, the inefficiency of conventional Learning Management Systems (LMS) is not merely a structural issue but a psychological one; "attentional overload" occurs when environmental demands exceed an individual's finite attentional capacity, hindering the deep processing required for career-critical skills\cite{Shan2023Impact}.

% Paragraph 2: Giải pháp đề xuất và công nghệ áp dụng (Tương đương đoạn 2 của bài mẫu)
By introducing an AI-powered Career Readiness Learning Management System, we aim to revolutionize the way educational institutions prepare students for the workforce. A core objective of aBridgeAI is to mitigate the risk of "AI as Autopilot," a failure mode where generative tools provide complete solutions that reduce "germane load"—the productive mental work needed to build mastery—and thus prevent actual learning. Instead of doing the thinking for the student, aBridgeAI utilizes Large Language Models (LLMs) to provide "fading support" scaffolding, which transitions learners from fully solved worked examples to independent practice as their competence grows. The system is designed to force "System-2 level engagement" by prompting students to explain their reasoning in steps, identify underlying assumptions, or solve problems using alternative methods. To further respect the constraints of human working memory, the platform employs "segmenting," breaking high-intrinsic-load curriculum into bite-size, learner-controlled segments. This ensures that the learner has sufficient time to organize and integrate information from one portion of the lesson before the next segment begins, preventing cognitive overload during the acquisition of complex technical skills\cite{studiomg_cognitive_ai, mayer2003}.

% Paragraph 3: Tác động vĩ mô và tính cấp thiết (Tương đương đoạn 3 của bài mẫu)
Moreover, the development of such a system addresses the pressing need for automation and scalable assessment in the educational sector while actively countering "digital dementia." This term refers to the cognitive decline and memory loss caused by overreliance on digital devices to store information that would otherwise be held in long-term memory. aBridgeAI combats this by integrating "active recall" and hybrid spacing algorithms that identify the optimal intervals for review just before information is likely to be forgotten. As students gain expertise, the system dynamically fades guidance to avoid the "expertise reversal effect," where redundant instructional support becomes an "extraneous load" that interferes with the schemas already constructed by the learner. This automated and context-aware approach allows educators to track authentic progress and provide personalized remediation, ensuring that mastery is internalized rather than offloaded to a tool. Ultimately, aBridgeAI answers the critical need to guide learners from being merely "Book Smart" to becoming "Job-Ready" candidates whose professional competencies are automated and resilient\cite{Renkl2003, wozniak1998sm2}.